\scanpage{163}
\mspage{37}
% archive sheet number: 162

$C$ et $F$, $TC$ et $TF$ sont clos à g et resp. à dr.). L'axiome
M5 est la saturation \deleted{\unsure{non}} modérée de $W$.

4.\deleted{\unsure{3}}4. Voici l'idée pour construire $\pm$ ``toutes'' les
\deleted{\unsure{couples}} structures de Quillen closes
$(C,TC,F,TF)$ \unsure{associée|associées} à un localiseur donné,
$W\subset M$. On suppose que $M$ satisfait aux
conditions du § précédent, concernant l'existence
de $\underrightarrow{\lim}$ filtrantes, d'une filtration \unsure{cardinale|ordinale},
l'accessibilité de $M$ \deleted{\unsure{(pour)}} \inserted{\unsure{celle|avec}} des éléments de
$\mathrm{Filt}^{\pi_0}$). On part d'une partie petite
\[
C_0\subset\mathrm{Fl}\,M
\]
telle que
\[
C_{0*}\subset W,\tag*{(\unsure{4|5}.4.1)}
\]
et on pose
\[
\left\{\begin{aligned}
&C=\widetilde C_0,\qquad TC=C\cap W.\\
&TF=C_*=C_{0*},\qquad F=(TC)_*.
\end{aligned}\right.\tag*{(4.4.2.)}
\]
\deleted{\unsure{On}}

Il y a \deleted{deux} \inserted{un} problème\unsure{s} seulement, vu
que l'on sait déjà que $(C,TF)$ est un
couple de Quillen par le théorème de factorisation
du § 1 (p. 10) : Il faut la factorisation pour
le couple $(TC,F=(TC)_*)$. À ce moment là,
les conditions préliminaires \inserted{(4.3.1), et (4.3.2) de prop.}
\deleted{\unsure{de}} 4.3. sont
satisfaites
\[
TC\subset C,\quad TF\subset F,\quad W\supset TC,TF\deleted{\unsure{\subset W}},
\]
ainsi que la condition a \quad $C\cap W\subset\widetilde{TC}$, puisque
$C\cap W=TC$. Donc on trouve, par l'équivalence
des conditions a, a', c) et vu que $TC$ et $TF$ sont clos
\[
\unsure{TF|TP}\unsure{\supseteq|=} F\cap W,\qquad W=TF\circ TC
\]
